\section{Defective cells and an insertion comparison}\label{sec:defects}

We combine occupation, isolation, and trajectory estimates to control
an arbitrary specified set of unusable cells. The final comparison
of this section transfers that control from a probability measure
on closed paths to the finite measure obtained by opening one end.

Fix
$\Delta_A=1000(2R_A+2R+20)^3$.
In this section all symmetric path tests use the actual inner
endpoints. A cell is \emph{defective} if its end count is below
$2b_-m$ or above $b_+m/2$, or if at least $\varepsilon/4$ of its
slots fail eligibility with score threshold $C_0-\Delta_A$
in place of $C_0$. Thus every unusable cell is defective.

\begin{lemma}[Scores at separated anchors]\label{lem:anchor-scores}
Let $U$ be a cell union and let $\mathcal D$ be any collection of
anchors mutually separated by more than $r$, all within distance
$Ds$ of $U$. There is a coarse-cell cover $S$ of the
$(D+R_A)s$ neighborhood of $U$, with
$|S|\leq C(D+R_A+2)^3|U|$, such that, pathwise,
\begin{equation}
 \sum_{a\in\mathcal D}\mathcal A^-(a;\omega)
 \leq C\int_0^{\delta/2}(w(t)+r^{-3})n_S(\omega(t))\,dt.
 \tag{21}\label{eq:21}
\end{equation}
The analogous inequality holds for the end scores after time
reversal. The assertion holds also when the anchors are chosen
from the paths by an arbitrary measurable rule, provided their
separation and location hold for each realization.
There is $B_0=C(1+K^2/\zeta^3)$ such that the exponential moment
of the right side of (21), divided by $B_0$, is at most
$\exp(C_v\rho|S|)$ under the closed measure.
\end{lemma}

\begin{proof}
Lift the anchors periodically to $\mathbb R^3$. Their distinct
lifts remain separated by more than $r$ once $L>2r$.
For a point $y$ and $a_t=\sqrt{1+t}$, packing disjoint balls
of radius $r/2$ bounds the number of anchors in a ball of radius
$u$ by $C(1+u/r)^3$. Summing in shells of width
$\max(r,a_t)$, with the Gaussian weight, gives
\[
 \sum_{a\in\mathcal D}
 e^{-d(y,a)^2/[1000(1+t)]}
 \leq C\bigl(1+(1+t)^{3/2}r^{-3}\bigr).
\]
For the torus expression, select a nearest lift of each anchor;
it is bounded by the full lifted Gaussian sum just estimated.
Only $y$ in $S$ can contribute after the distance indicator
in (15), because the anchors are within $Ds$ of $U$.
Multiplication by $w(t)$ now proves (21), since
$w(t)(1+t)^{3/2}\leq C$. This proof is deterministic for every
set of separated anchors and therefore permits adaptive choices.

The integral of the scalar weight on the right of (21) is at
most $C(1+\delta r^{-3})\leq B_0$.
Divide by $B_0$ and regard the result as an average of the
variables $n_S(\omega(t))$, padding with a zero variable if
the total averaging mass is less than one. Jensen's inequality
for the exponential, stationarity of the one-time closed law,
and (11) at $\lambda=1/2$ give
\[
 \mathbb E_{\mathrm{cl}}\exp\left[
 B_0^{-1}C\int_0^{\delta/2}(w(t)+r^{-3})n_S(\omega(t))\,dt
 \right]
 \leq\mathbb E_Pe^{n_S}
 \leq\mathbb E_Pe^{n_{\phi_S}}
 \leq e^{C_v\rho|S|}.
\]
Time reversal gives the other half of the assertion.
\end{proof}

\begin{proposition}[Specified defective cells]\label{prop:defective-cells}
Choose $D$ and $\varepsilon$ so that
$C_v\sum_{j\geq0}D_j^3e^{-\kappa D_j^2}
\leq10^{-4}b_-\varepsilon$ in (17). Then choose $\zeta>0$
sufficiently small in terms of these fixed constants, followed
by $C_0>\Delta_A$ sufficiently large. For small enough density,
there is $q_{\mathrm{cl}}=q_{\mathrm{cl}}(\rho)\to0$ as
$\rho\downarrow0$ such that every collection $\mathcal U$
of specified cells satisfies
\begin{equation}
 P_{\mathrm{cl}}(\text{every cell of }\mathcal U
                         \text{ is defective})
 \leq q_{\mathrm{cl}}^{|\mathcal U|}.
 \tag{22}\label{eq:22}
\end{equation}
The bound and the convergence of $q_{\mathrm{cl}}$ are uniform
in the allowed temperatures and sufficiently large volumes.
Increasing $C_0$ preserves the statement.
\end{proposition}

\begin{proof}
Let $u=|\mathcal U|\geq1$. On the event in question, some
subcollection of at least $u/3$ cells shares one of three
reasons: low count, high count, or at least
$\varepsilon/4$ ineligible slots with counts in
$[2b_-m,b_+m/2]$. For the first two reasons (12) applies
to the union of that subcollection.

Fix a subcollection of $u'$ cells of the third kind, and let
$U'$ be its union. It contains at least
$\varepsilon b_-mu'/2$ ineligible slots. Assign to each such
slot the first failed test in the following ordered list:
trajectory failure; end nonisolation; start nonisolation;
beginning score failure; end score failure. For each test
after the first, the slot is macro-conforming; for each score
test, its respective anchor is isolated. At least one reason
therefore has at least
$k'=\varepsilon b_-mu'/20$ slots. Using this slightly smaller
threshold avoids any integer-rounding convention.

For trajectories, exponential Markov and (17) give probability
at most
\[
 \exp\{[10^{-4}b_-\varepsilon+o_\rho(1)]\rho|U'|-k'\},
\]
which is at most $e^{-ck'}$ for sufficiently small density,
since $\rho|U'|\leq2mu'$.
For end nonisolation use (13) with $S=U'$.
For start nonisolation, macro conformity places all those starts
within distance $Ds$ of $U'$. A coarse cover of that neighborhood
has volume at most $C(D+2)^3|U'|$, so (13) applies to it at
time zero. Fix $\zeta$ so small that in both applications
$C_v\zeta^3\rho|S|\leq k'/2$.
Their probabilities are then at most $e^{-\lambda_0 k'/2}$.

For either score reason the failed anchors are mutually separated
by more than $r$. The end anchors lie in $U'$, and the start
anchors lie within distance $Ds$ of $U'$. If at least $k'$
of them have score greater than $C_0-\Delta_A$, their score
sum exceeds $(C_0-\Delta_A)k'$. Apply
Lemma~\ref{lem:anchor-scores} and exponential Markov. The
probability is at most
\[
 \exp\{C_v\rho|S|-(C_0-\Delta_A)k'/B_0\}.
\]
The ratio $\rho|S|/k'$ is bounded in terms of the already
fixed constants. Choose $C_0$ large enough that this last
bound is at most $e^{-k'}$.

For small density, $\lambda_0=\zeta\rho^{1/6}\leq1$.
All the estimates for a fixed subcollection are thus bounded
by a fixed constant times $e^{-c\lambda_0mu'}$, including
the two count reasons. There are at most $2^u$ choices of
a subcollection for each reason and five failed-test choices.
Since $u'\geq u/3$, the resulting bound has the form
$C_1^u e^{-c\lambda_0mu/3}$ after increasing $C_1$.
Now
\[
 \lambda_0m\geq\tfrac12\zeta K^3\rho^{-1/3}\longrightarrow\infty.
\]
Absorbing $C_1^u$ proves (22), for example with
$q_{\mathrm{cl}}=C_1e^{-c\lambda_0m/3}$ after replacing
$m$ by its stated density-dependent lower bound. Each preceding
bound was uniform in temperature and volume, as required.
The empty collection has probability one and satisfies (22).
\end{proof}

We have obtained a joint defect bound under the closed probability
law. The second-moment comparison will instead produce a finite
measure with one endpoint opened. The following transfer retains both
the small defect parameter and the cell-pair mass factor $M^{-2}$.

For cells $a,b$, define the positive measure $\mu_{ab}$ as
follows. In (14), sum over a slot $i$, restrict its outer endpoint
$Y_i$ to cell $b$, and replace its inner endpoint by a point
$z\in a$ integrated with Lebesgue measure. All other inner
ends remain $Y_j$. Multiply by $1/(NV)$, retain normalization
by $Z_N$, and integrate all the other variables as before.
The slot index may be retained as part of the measure space.
An event on the actual slab paths is \emph{symmetric} when
it is invariant under simultaneous relabeling of all slots;
it must not single out the opened slot.

\begin{proposition}[Transfer to an open insertion]\label{prop:insertion-comparison}
For any measurable symmetric slab event $\mathcal E$,
\begin{equation}
 \mu_{ab}(\mathcal E)
 \leq\frac C{M^2}P_{\mathrm{cl}}(\mathcal E)^{1/4}.
 \tag{23}\label{eq:23}
\end{equation}
In particular, the insertion mass is at most $C/M^2$ and
the joint bound (22) transfers to this measure with exponent
$q_{\mathrm{cl}}^{|\mathcal U|/4}$ and prefactor $C/M^2$.
All tests in this application use the actual inner endpoints.
\end{proposition}

\begin{proof}
Two applications of Cauchy--Schwarz give the fourth root: the first
separates the two insertion factors, and the second bounds a local
particle count on the event $\mathcal E$ by its second moment.
Let $K_{\mathcal E}(E,S)$ be the tested slab kernel obtained
by integrating $W(S,E;d\omega)\mathbf1_{\mathcal E}$.
It need not be self-adjoint, but it is pointwise nonnegative,
is bounded above by the untested slab kernel, and commutes
with simultaneous coordinate permutations. Put
$U_{\mathcal E}=K_{\mathcal E}e^{-(\beta/2-\delta)H}$.
For a one-particle function $f$, let $a(f)$ denote annihilation,
mapping the $N$-particle symmetric space to the $(N-1)$-particle
one. The usual normalization gives
$a^*(\mathbf1_b)a(\mathbf1_a)
=\sum_i(|\mathbf1_b\rangle\langle\mathbf1_a|)_i$.
The definition of $\mu_{ab}$ and semigroup cyclicity therefore give
\[
 \mu_{ab}(\mathcal E)=\frac1{NVZ_N}
 \operatorname{Tr}_{\mathrm{sym}}
 e^{-\beta H/2}a^*(\mathbf1_b)a(\mathbf1_a)U_{\mathcal E}.
\]
Apply Hilbert--Schmidt Cauchy--Schwarz to
$a(\mathbf1_b)e^{-\beta H/2}$ and
$a(\mathbf1_a)U_{\mathcal E}$. Since a cell has volume $s^3$,
the one-body Cauchy--Schwarz inequality gives
$a^*(\mathbf1_a)a(\mathbf1_a)\leq s^3n_a$, and similarly
for $b$. Translation invariance yields
$\operatorname{Tr}e^{-\beta H}n_b=Z_Nm$.
Thus the first squared Hilbert--Schmidt norm is at most
$s^3Z_Nm$.

For the second norm use
$\operatorname{Tr}U_{\mathcal E}^*n_aU_{\mathcal E}$.
This is an integral of products of nonnegative kernels.
Precisely, on labeled coordinates one may insert the projector
and use the kernel of $U_{\mathcal E}P_{\mathrm{sym}}$; it
has the same trace on its symmetric range and obeys
\[
 0\leq U_{\mathcal E}P_{\mathrm{sym}}
 \leq e^{-\beta H/2}P_{\mathrm{sym}}
 \quad\hbox{pointwise in kernels}.
\]
Replace only the adjoint kernel in that trace by the untested
one. This increases the integral and, by cyclicity, gives
\[
 \operatorname{Tr}_{\mathrm{sym}}U_{\mathcal E}^*n_aU_{\mathcal E}
 \leq\operatorname{Tr}_{\mathrm{sym}}
 e^{-(\beta-\delta)H}n_aK_{\mathcal E}
 =Z_N\mathbb E_{\mathrm{cl}}[n_a\mathbf1_{\mathcal E}].
\]
Here $n_a$ counts the actual inner endpoints. Its second
moment in the closed measure is at most $Cm^2$ by (11),
so the last expression is at most
$CZ_NmP_{\mathrm{cl}}(\mathcal E)^{1/2}$.
The second squared Hilbert--Schmidt norm is consequently at
most $Cs^3Z_NmP_{\mathrm{cl}}(\mathcal E)^{1/2}$.
Taking square roots and dividing by $NVZ_N$ gives
$Cs^3m/(NV)$ times the fourth root of the probability.
Since $m=N/M$ and $s^3=V/M$, this factor is $C/M^2$,
which proves (23).
\end{proof}
